\chapter{Evaluationsmethodik}\label{ch:methodology}

Dieses Kapitel beschreibt, wie gemessen wird; die konkreten Ergebnisse folgen in
Kapitel~\ref{ch:results}.

\section{Drei Messreihen}\label{sec:series}

Die Architektur sieht drei Messreihen vor, die alle vom \texttt{messung\_\allowbreak driver} über die
\texttt{Experiment\allowbreak Driver}-Bibliothek ausgeführt werden, konfiguriert entweder als drei feste
Reihen oder über eine externe \texttt{messreihen.xml}.

\begin{description}
  \item[Reihe A --- PRT-ART vs.\ Stand der Technik.] Vergleicht den Prüfling gegen die
  dokumentierten SOTA-Algorithmen. Zwei Modi: \emph{A\_defined} (schnelle Verifikation
  gegen die 8 Rang-1-Profile) und \emph{A\_full} (alle 30 SOTA-Profile; der Standard,
  gemäß der Direktive, so vollständig wie möglich zu messen).
  \item[Reihe B --- Cache-Engine-Permutationen.] Die reine ``Stand der Technik''-Vergleichsbasis:
  alle SOTA-Profile ohne den PRT-ART-Prüfling.
  \item[Reihe C --- Merge alt/neu.] Konkrete Merge-Punkte zwischen PRT-ART-Bausteinen und
  Cache-Engine-Permutationen, z.\,B.\ PRT-ART OLC + tcmalloc + ART-Page; PRT-ART
  4+2-Pool + HOT-Page; PRT-ART Path-Prefetch + B\textsuperscript{2}-Baum-Page; PRT-ART
  MultiLevel-Layout + Wormhole-Page.
\end{description}

Jede dieser drei Reihen wird in drei Granularitäten erhoben (Aufgabenstellung-Methodik): Das
\emph{Micro-Benchmarking} vergleicht jeden einzelnen Entwurfsbestandteil (Achse) isoliert
über das gemeinsame Interface seiner Kategorie gegen Lastprofile; das \emph{Makro-Benchmarking}
misst die einzelnen \texttt{std::map}-Interface-Operationen (Einfügen, Punkt- und
Bereichssuche, Aktualisieren, Löschen) gegen Beispiel-Lasten; und das \emph{Gesamt-Benchmarking}
fährt die etablierten YCSB-Lastprofile über jede Rekombination der Bausteine. Aus dem Ranking
über diese drei Granularitäten ergibt sich, welche Entscheidung an welcher Stelle für welches
Lastmuster vorteilhaft ist.

\section{ExperimentDriver-Phasen}\label{sec:experimentdriver}

Der \texttt{ExperimentDriver} durchläuft sieben Phasen je Reihe: (1)~\textbf{Enumerate}
(XML-Configs parsen und die Permutationsliste aufbauen); (2)~\textbf{Codegen} (ein
\texttt{comdare\_perm\_<fp>.cpp} je Permutation, plus
\texttt{generate\_module\_from\_profile} für die SOTA-Profile); (3)~\textbf{Compile};
(4)~\textbf{Load} (LoadLibrary/dlopen aller Module); (5)~\textbf{Execute}
(\texttt{create\_instance} + profilbewusstes \texttt{run\_workload}); (6)~\textbf{Measure}
(die binären Mess-Records aggregieren); (7)~\textbf{Persist} (\texttt{measurements.csv} +
\texttt{.json}). Zwei opt-in-Phasen kompilieren fehlende Module hot und verifizieren den
ABI-Vertrag je Modul.

\section{Hypothesen}\label{sec:hypotheses}

\begin{description}
  \item[H1 (Page-Type-Kosten)] Die mittleren Zugriffskosten je Page-Type-Variante
  unterscheiden sich systematisch nach Workload: dichte Page-Typen gewinnen
  erwartungsgemäß unter Zipf-Verteilung (YCSB-A), bereichsoptimierte Pages unter
  Bereichs-Scans (YCSB-E).
  \item[H2 (Code-Qualität pro Quelle)] Der vom Original-Quell-Repository geerbte
  Code-Qualitäts-Score korreliert messbar mit dem erreichten Durchsatz.
  \item[H3 (Inline- vs.\ External- vs.\ Chain-Verteilung)] Die beobachtete Verteilung der drei
  ValueHandle-Varianten korreliert mit der Page-Dichte: dichte Pages bevorzugen Inline,
  spärliche Pages External und PRT-ART-Heuristiken ChainRef.
\end{description}

\section{Workloads und Datensätze}\label{sec:workloads}

Phase~5 leitet den YCSB-Workload je Permutation aus dem Traversal-Tag des Profils ab
(Tabelle~\ref{tab:workload-routing}); Module ohne Profil verwenden den globalen
Default-Workload.

\begin{table}[h]
  \centering
  \caption{Profil-bewusstes Workload-Routing}\label{tab:workload-routing}
  \begin{tabular}{lll}
    \toprule
    \textbf{Traversal-Tag} & \textbf{YCSB-Workload} & \textbf{Charakteristik} \\
    \midrule
    \texttt{RANGE\_SCAN} / \texttt{RANGE\_HOT\_PATH} & YCSB-E & Bereichsabfragen \\
    \texttt{HOT\_PATH} / \texttt{PRTART\_HOT\_PATH} & YCSB-A & 50/50 Lesen/Aktualisieren, Zipf \\
    \texttt{STANDARD} (Default) & YCSB-C & nur Lesen \\
    \bottomrule
  \end{tabular}
\end{table}

Als Datengrundlage dienen reale String-Datensätze unterschiedlicher Charakteristik
(Tabelle~\ref{tab:datasets}); jeder trägt eine Reproduzierbarkeits-Akte aus Quelle, Prüfsumme,
Zeilenzahl, Vorverarbeitung und Seed-Regel für die synthetischen Miss-, Update- und
Delete-Schlüssel.

\begin{table}[h]
  \centering
  \caption{Reale Datensätze und ihre Charakteristik}\label{tab:datasets}
  \begin{tabular}{ll}
    \toprule
    \textbf{Datensatz} & \textbf{Charakteristik} \\
    \midrule
    \texttt{url} & sehr lange Schlüssel, lange gemeinsame Präfixe, mittleres Alphabet \\
    \texttt{dna} & kleines Alphabet, hohe Wiederholung, tiefe Präfix-Entropie \\
    \texttt{protein} & biologische Sequenzen, kleines bis mittleres Alphabet \\
    \texttt{xml} & strukturierte Strings mit wiederkehrenden Präfixen \\
    \texttt{tpcds-id} & datenbanknahe IDs, sortiert und unsortiert \\
    \texttt{trec-terms} & Wörterbuch-/Suchindex-Terme mit definierten Miss-Schlüsseln \\
    \bottomrule
  \end{tabular}
\end{table}

\section{Versuchsplattformen}\label{sec:platforms}

Die Messungen zielen auf eine lokale Workstation zur Entwicklungs-Verifikation und den
TU-Dresden-ZIH-Cluster (Barnard CPU, Capella GPU) für die vollständigen Sweeps, plus die
Rang-3-Plattformen (Ryzen~9~9950X3D, Intel~i9~14900KS und eine HBM-Plattform). Ein
\texttt{IPlatformProbe} (Phase~1) meldet je Plattform die lauffähige Teilmenge der Achsen
(z.\,B.\ AVX-512 nur dort, wo verfügbar, HBM-aware nur auf HBM-Hardware). Je Plattform werden
ein Scalar-, ein AVX2- und --- auf Barnard --- ein AVX-512-Build vorbereitet; auf Hybrid-CPUs
werden P- und E-Cores getrennt vermessen (eigene Zähler-Domänen
\texttt{cpu\_core}/\texttt{cpu\_atom}) statt zusammenaggregiert.

\smallskip\noindent\emph{Experiment-Betriebssysteme.} Auf den Produktions-Zielmaschinen wird jede Messung
unter zwei Betriebssystem-Regimes erhoben, um \emph{Produktions-Treue} von \emph{Mess-Tiefe} zu trennen. Ein
\emph{immutables} Betriebssystem (Talos) spiegelt den Produktiv-Einsatz und erzwingt boot- bzw.\
übersetzungsstatische Cache-Einstellungen (Abschnitt~\ref{ssec:aware-oblivious}); ein \emph{root-Linux mit
vollem Hardware-Zähler-Zugriff} (\texttt{perf}/MSR) macht hingegen die zählerbasierten Mess-Kategorien ---
Cache-, dTLB- und Branch-Misses sowie IPC/CPI (Abschnitt~\ref{ssec:sota-meas-bridge}) --- und die
laufzeit-konfigurierbaren cache-nahen Achsen-Parameter (Prefetch-Distanz, Hardware-Prefetcher) überhaupt erst
zugänglich. Die zeit- und observer-basierten Kategorien (Wall-Clock, HdrHistogramm, Per-Achsen-Observer)
laufen unter beiden Regimes identisch und sichern die Vergleichbarkeit; die PMC-Kategorien werden nur unter dem
privilegierten Regime erhoben und je Plattform mit Betriebssystem, Kernel und Zähler-Domäne protokolliert.

\section{Fairness und Reproduzierbarkeit}\label{sec:fairness}

Faire Vergleichbarkeit ist über feste Regeln gesichert. Jeder Vergleich trennt einen
gemeinsamen Minimalmodus (Common-Denominator: externe Werte-Handles, keine
PRT-ART-Spezialpfade) vom PRT-ART-Native-Modus (Inline-Umschaltung, Cache-Engine,
Seitentyp-Scheduler). Alle Verfahren erhalten identische, normalisierte Byte-Schlüssel sowie
identische Miss-, Update- und Delete-Mengen; Warmup, Randomisierung und Thread-Pinning werden
je Plattform und Workload protokolliert. Fehlende Fähigkeiten einer Baseline --- etwa
Bereichs-Scans einer Hash-Tabelle --- werden als N/A ausgewiesen statt über Hilfsstrukturen
verdeckt emuliert; für jede Fremdbibliothek werden Compiler, Flags, ISA-Pfad, Allokator und
Commit-Hash mitgeloggt. Für die Reproduzierbarkeit wird jede Konfiguration in mehreren
getrennten Läufen vermessen, deren Streuung Teil der Aussage bleibt: Perzentile werden über
HDR-Histogramme bestimmt statt gemittelt.

\section{Permutations-Explosion und Reduktion}\label{sec:explosion}

Mit der N-Phasen-Erweiterung von 11 auf 19 Achsen plus $\sim$57 Sub-Achsen wächst der Raum von
$\sim$5,5 Milliarden auf weit mehr als $10^{11}$ Permutationen. Drei Mechanismen dämpfen die
Explosion: (1)~der \texttt{<expected\_workload>}-Profil-Filter routet nur kompatible
Workloads; (2)~\texttt{MessreihenMode::Defined} führt nur die deklarierten Tupel aus
($\sim$100--1000 je Reihe), die in das Manuskript einfließen; (3)~\texttt{Full} (und das
realistischere \texttt{Full-Sampled}, z.\,B.\ jede 1000.\ Permutation) läuft nur auf dem
ZIH-Cluster, unter Beachtung der vom Plattform-Probe gemeldeten Achsen-Kompatibilität.
